\documentclass[11pt,letterpaper]{article}

\usepackage[spanish,es-noshorthands,es-tabla]{babel}
\usepackage{fontspec}
\setmainfont{Source Serif Pro}
\usepackage{microtype}
\usepackage{geometry}
\usepackage{enumitem}
\usepackage{xcolor}
\usepackage{csquotes}
\usepackage{float}
\usepackage{booktabs}
\usepackage{array}
\usepackage{tikz}
\usepackage{pgfplots}
\pgfplotsset{compat=1.18}
\usetikzlibrary{arrows.meta,positioning,fit,backgrounds,calc}
\usepackage{xurl}
\usepackage[hidelinks]{hyperref}
\usepackage[backend=biber,style=apa,sorting=nyt,maxcitenames=2]{biblatex}
\DeclareLanguageMapping{spanish}{spanish-apa}

\geometry{margin=2.3cm}
\hbadness=10000
\setlength{\emergencystretch}{1em}
\setlength{\parindent}{0pt}
\setlength{\parskip}{0.5em}
\setlist{nosep}
\addbibresource{referencias.bib}
\AtBeginBibliography{\sloppy\small}
\setlength{\bibitemsep}{0.3em}

% Colores: azul petróleo para el presente, ámbar para las proyecciones,
% azul oscuro para rótulos.
\definecolor{cHoy}{RGB}{31,111,139}
\definecolor{cFut}{RGB}{214,137,16}
\definecolor{cEje}{RGB}{26,58,102}
\definecolor{cGris}{RGB}{120,116,110}

\pgfplotsset{
  foro/.style={
    /pgf/number format/use comma,
    /pgf/number format/1000 sep={},
    tick label style={font=\small}, axis line style={cGris}, tick style={cGris},
    ylabel style={font=\small, text=cEje}, xlabel style={font=\small, text=cEje},
    ymajorgrids, grid style={cGris!25},
  },
}

\begin{document}

% --- Portada -----------------------------------------------------------------
\begin{titlepage}
\centering
\vspace*{3.2cm}
{\large\scshape Gestión del Desarrollo Tecnológico:\\
Estrategias y Análisis de Portfolio\par}
\vspace{0.5cm}
{\normalsize Módulo 9. Foro 9.1: Relación entre edad y tecnología\par}
\vspace{2.6cm}
{\color{cEje}\rule{0.55\textwidth}{0.8pt}\par}
\vspace{0.7cm}
{\LARGE\bfseries La vejez es una frontera móvil\par}
\vspace{0.5cm}
{\large Envejecimiento, tecnología y el mercado\\
que la industria todavía no mira\par}
\vspace{0.7cm}
{\color{cEje}\rule{0.55\textwidth}{0.8pt}\par}
\vfill
{\large Carlos Luis Rojas Aragonés\par}
\vspace{0.3cm}
{\normalsize Chief Technology Officer Program\par}
\vspace{0.3cm}
{\normalsize 2 de noviembre de 2026\par}
\vspace*{1.5cm}
\end{titlepage}

% --- Cuerpo ------------------------------------------------------------------
\section*{1. ¿Cuándo empieza la vejez?}

Creo que no es un umbral objetivo, sino una variable que depende del contexto:
la esperanza de vida, la distribución de edades de la población y hasta la
ley de cada país. La OMS habla de personas de 60 años o más
\parencite{who2024}, mientras que en Costa Rica la Ley Integral para la Persona
Adulta Mayor fija el corte en 65 años \parencite{ley7935}. Son convenciones, y
las convenciones envejecen.

Los demógrafos Sanderson y Scherbov lo plantean de una forma que me parece muy
útil: en lugar de contar los años vividos, contar los que quedan. Proponen
considerar \enquote{vieja} a una persona cuando su esperanza de vida restante
baja de unos 15 años \parencite{sanderson2010}. Con esa regla, la vejez no
llega a la misma edad en todos los países ni en todas las épocas: se corre
hacia adelante a medida que la medicina y la tecnología mejoran. En Costa Rica,
la esperanza de vida al nacer pasó de 76,95 años en 1990 a 80,91 en 2024, y el
INEC proyecta 84,27 en 2050 y 89,57 en 2100 \parencite{ramirez2024}
(Figura~\ref{fig:esperanza}). Si eso se cumple, vamos a empezar a considerar
jóvenes a muchas personas que hoy llamamos mayores.

\begin{figure}[H]
\centering
\begin{tikzpicture}
\begin{axis}[
  foro, width=11cm, height=6.2cm,
  xmin=1980, xmax=2110, ymin=70, ymax=95,
  xtick={1990,2024,2050,2100},
  ylabel={Años},
]
  \addplot[cFut, line width=1.4pt, dashed, mark=*, mark size=2.6pt, mark options={solid}]
    coordinates {(2024,80.91) (2050,84.27) (2100,89.57)};
  \addplot[cHoy, line width=1.4pt, mark=*, mark size=2.6pt]
    coordinates {(1990,76.95) (2024,80.91)};
  \foreach \x/\y/\t in {1990/76.95/{76,95}, 2024/80.91/{80,91},
                        2050/84.27/{84,27}, 2100/89.57/{89,57}} {
    \edef\temp{\noexpand\node[font=\noexpand\small\noexpand\bfseries, anchor=south, yshift=3pt]
      at (axis cs:\x,\y) {\t};}
    \temp
  }
\end{axis}
\end{tikzpicture}
\caption{Esperanza de vida al nacer en Costa Rica: valores estimados (línea
continua) y proyectados (línea discontinua). Elaboración propia con datos del
INEC citados por \Textcite{ramirez2024}.}
\label{fig:esperanza}
\end{figure}

\section*{2. El efecto social: trabajo, salud y tecnología}

El cambio es global. La OMS estima que para 2030 una de cada seis personas en
el mundo tendrá 60 años o más, y que ese grupo pasará de 1000 millones en 2020
a 2100 millones en 2050; los mayores de 80 se triplicarán hasta 426 millones
\parencite{who2024}. Según la ONU, a mediados de la década de 2030 habrá más
personas de 80 años o más que bebés, y a finales de la de 2070 los mayores de
65 superarán a los menores de 18 \parencite{un2024}.

Costa Rica va más rápido que el promedio. En 2024, 11 de cada 100 habitantes
tenían 65 años o más; en 2050 serán 25. La relación entre mayores de 65 y
personas en edad de trabajar subirá de 16,0 a 39,2 por cada 100, y la población
total dejará de crecer en 2044 \parencite{ramirez2024} (Figura~\ref{fig:cr}).

\begin{figure}[H]
\centering
\begin{tikzpicture}
\begin{axis}[
  foro, ybar, bar width=30pt, width=11cm, height=6.4cm,
  enlarge x limits=0.4, ymin=0, ymax=48,
  symbolic x coords={pob,dep}, xtick=data,
  xticklabels={{Por cada 100 habitantes},{Por cada 100 personas\\en edad de trabajar}},
  xticklabel style={align=center},
  ylabel={Personas de 65 años o más},
  nodes near coords, nodes near coords style={font=\small\bfseries},
  /pgf/number format/fixed, /pgf/number format/precision=1,
  legend style={at={(0.03,0.97)}, anchor=north west, draw=none,
    fill=none, font=\small, legend cell align=left},
  legend image code/.code={\draw[#1, fill] (0cm,-0.1cm) rectangle (0.3cm,0.12cm);},
]
  \addplot[draw=cHoy, fill=cHoy!80] coordinates {(pob,11) (dep,16.0)};
  \addplot[draw=cFut, fill=cFut!80] coordinates {(pob,25) (dep,39.2)};
  \legend{2024, 2050 (proyección)}
\end{axis}
\end{tikzpicture}
\caption{Peso de la población de 65 años o más en Costa Rica. Elaboración
propia con datos del INEC citados por \Textcite{ramirez2024}.}
\label{fig:cr}
\end{figure}

\textbf{En lo laboral}, la lectura pesimista es que habrá menos personas
cotizando por cada pensionado. Pero si la gente envejece mejor, también va a
estar en mejores condiciones de seguir trabajando, emprendiendo o cambiando de
campo mucho después de los 65. Por eso Sanderson y Scherbov sostienen que la
edad de retiro no debería ser fija, sino moverse con la esperanza de vida
\parencite{sanderson2010}.

\textbf{En lo médico}, vivir más significa convivir más años con enfermedades
crónicas, y eso empuja la demanda de servicios de salud y de cuidados
\parencite{who2024}. La medicina está avanzando bastante, y la presión
demográfica es justamente uno de los motores de la telemedicina, el monitoreo
remoto y los dispositivos que se usan en el cuerpo.

\textbf{En lo tecnológico}, el envejecimiento mueve dinero. AARP y Economist
Impact estimaron que en 2020 las personas de 50 años o más, el 24\,\% de la
población mundial, generaron el 34\,\% del PIB global (unos 45 billones de
dólares), y proyectan que en 2050 llegarán al 39\,\% (118 billones)
\parencite{aarp2022} (Figura~\ref{fig:longevidad}).

\begin{figure}[H]
\centering
\begin{tikzpicture}
\begin{axis}[
  foro, ybar, bar width=40pt, width=10cm, height=6.2cm,
  enlarge x limits=0.5, ymin=0, ymax=140,
  symbolic x coords={2020,2050}, xtick={2020,2050},
  xticklabels={2020, 2050 (proyección)},
  ylabel={Billones de dólares},
  nodes near coords, nodes near coords style={font=\small\bfseries},
  bar shift=0pt,
]
  \addplot[draw=cHoy, fill=cHoy!80, bar shift=0pt] coordinates {(2020,45)};
  \addplot[draw=cFut, fill=cFut!80, bar shift=0pt] coordinates {(2050,118)};
  \node[font=\footnotesize\bfseries, text=white, align=center, anchor=north]
    at (axis cs:2020,40) {34\,\%\\del PIB};
  \node[font=\footnotesize\bfseries, text=white, align=center, anchor=north]
    at (axis cs:2050,110) {39\,\%\\del PIB};
\end{axis}
\end{tikzpicture}
\caption{Aporte de las personas de 50 años o más al PIB mundial. Elaboración
propia con datos de \Textcite{aarp2022}.}
\label{fig:longevidad}
\end{figure}

\section*{3. Retos y oportunidades para el sector tecnológico}

Como vimos en uno de los videos del curso, se abre un mercado nuevo: el de las
personas que ya se retiraron, o que están por hacerlo, y que pueden convertirse
fácilmente en un público objetivo para empresas que tradicionalmente apuntaban
a poblaciones cada vez más jóvenes. Joseph Coughlin, director del MIT AgeLab,
lo llama la \enquote{economía de la longevidad} y advierte que las empresas
ven el tamaño y la riqueza de ese mercado, pero lo atienden con una idea
anticuada de lo que es ser mayor: productos que gritan \enquote{para viejos} y
que nadie quiere usar \parencite{coughlin2017}.

Veo tres retos y tres oportunidades:

\begin{itemize}
  \item \textbf{Reto: dejar de diseñar para el estereotipo.} El usuario mayor
    ya está conectado: en Estados Unidos, el 61\,\% de los mayores de 65 tenía
    teléfono inteligente en 2021, frente al 96\,\% de los de 18 a 29
    \parencite{faverio2022}. La brecha existe, pero se cierra rápido; lo que
    falla es el diseño, no el interés.
  \item \textbf{Reto: accesibilidad como requisito, no como extra.} Vista,
    oído, motricidad fina y carga cognitiva cambian con la edad. Un producto
    accesible para alguien de 75 suele ser mejor para todos.
  \item \textbf{Reto: el sesgo de edad dentro de los propios equipos.} Es
    difícil diseñar bien para un público que no está representado en la sala.
  \item \textbf{Oportunidad: salud y cuidado.} Telemedicina, monitoreo remoto,
    detección de caídas y logística de medicamentos, como el modelo de
    farmacia que Coughlin cita como ejemplo \parencite{coughlin2017}.
  \item \textbf{Oportunidad: aprendizaje y trabajo a lo largo de la vida.}
    Carreras de 50 años o más necesitan plataformas para reentrenarse y
    cambiar de oficio.
  \item \textbf{Oportunidad: finanzas y vivienda.} Más años de retiro exigen
    productos para administrar el ahorro y para envejecer en casa.
\end{itemize}

Para quienes trabajamos en tecnología hay además una lectura personal: nuestras
propias carreras van a ser más largas, así que el aprendizaje continuo deja de
ser opcional.

Les dejo una pregunta: en sus empresas, ¿las personas mayores de 60 aparecen
en los perfiles de usuario y en las pruebas de producto, o siguen siendo un
segmento invisible?

\printbibliography[title={Referencias}]

\end{document}
